\section{Results and Discussion}
\label{sec:results}

\subsection{Main comparison}

The geometric reference scores only registered pairs (48\% of valid query--gallery pairs, containing all 5,722 positive pairs), whereas crop-based methods score all valid pairs; the two result sets use different candidate pools.
With \regonly{} (a constant on registered pairs, $-\infty$ otherwise), Rank-1 is 0.515, mAP 0.533, and DIR 0.000 at FAR 0.1 (zero because all registered pairs score equally). Adding mask comparison improves Rank-1 and mAP by 0.437 and 0.353, respectively, and raises DIR to 0.768 (Table~\ref{tab:main}). The highest crop-based DIR values are 0.429 (SuperGlue) and 0.404 (OSNet+ctx) (Section~\ref{sec:controls}), against chance 0.038.
No wall-clustered confidence intervals are reported.
At the validation-calibrated threshold, \ours{} exceeds its constant baseline by only 0.038 (0.582 vs.\ 0.544), compared with 0.112 at the oracle threshold ($F_1^{*} = 0.656$). This excess is smaller than LoFTR's 0.046 over its 0.305 baseline, suggesting that the contrast seen in ranking and DIR overstates the difference in thresholded decisions.
Registration gating may not help across visits; an appearance baseline restricted to registered pairs and the frame-gap baseline (ranking by the number of frames between photographs) remain to be evaluated.

\subsection{Registration rate (Section~\ref{sec:acquisition}) and failure modes}
\label{sec:coverage}
\label{sec:chamferablation}

The registration gate is label-informative: it retains all positives but only 36.8\% of negatives, which raises the prevalence from 0.180 to 0.373 and the constant-baseline $F_1$ from 0.305 to 0.544.

Registration failures depend on wall surface (Section~\ref{sec:cascade}); their extent on smooth plaster is unquantified. All 280 known queries have a true match in an adjacent frame (median gap 1), motivating the frame-gap baseline.

We define sharpness as the variance of the Laplacian of the image at quarter resolution.
A logistic regression on the test set suggests sharpness has a larger standardized association with registration success than frame gap (wall-clustered $z = 5.11$ vs.\ $-2.78$); with 11 clusters, however, the $p$-values overstate precision, and the sharpness threshold of 10, used to predict registration success, was chosen post hoc.
Registration, ECC, and crack-exclusion ablations were run on the test set and are not reported here; geometry settings should be fixed on validation before any final comparison.

\subsection{Synthetic revisit probe}
\label{sec:hybridbench}

Ten photographs from five walls were altered by a small common warp under condition A or B (Section~\ref{sec:protocol}), with one random seed.
Altered images were matched against real photographs. The warp is small (scale 1.08, rotation 6$^{\circ}$, out-of-plane tilt 4$^{\circ}$), whereas real test pairs reach 75$^{\circ}$ out-of-plane tilt at the 90th percentile (procedure in the code).
Table~\ref{tab:hybrid} compares \hybrid{} and OSNet+ctx.
For ranking and open-set identification (A/B), OSNet+ctx outperforms \hybrid{} (Rank-1 0.826/0.671 vs.\ 0.275/0.274; DIR 0.739/0.663 vs.\ 0.343/0.342; $F_1^{@v}$ similar at 0.601/0.594 vs.\ 0.569/0.574); assignment $F_1^{a}$ reverses narrowly (0.631/0.609 vs.\ 0.646/0.640), and no pair was scored by homography alone, so the reversal does not come from a homography-only path.
For RQ3, this single-seed probe provides no evidence of real multi-visit robustness, which remains untested (Section~\ref{sec:limitations}).

\begin{table}[t]

\centering

\scriptsize\setlength{\tabcolsep}{2pt}

 \caption{Synthetic revisit probe: \hybrid{} versus OSNet+ctx for A (growth) and B (photometric change), one seed, five walls. $n_{Q}$ number of queries; R@1 Rank-1 accuracy; mAP mean average precision; DIR detection and identification rate at FAR\,0.1; $F_{1}^{@v}$ validation-calibrated $F_1$; $F_{1}^{a}$ assignment $F_1$. \hybrid{} adds a crack-skeleton score to alignment-based comparison; \ours{} uses registration-gated Chamfer only. Bold: best method per column and condition.}

\label{tab:hybrid}

\begin{tabular}{llcccccc}

\toprule

Condition & Method & $n_{Q}$ & R@1 & mAP & DIR & $F_{1}^{@v}$ & $F_{1}^{a}$ \\

& & & & & {\scriptsize @FAR$_{0.1}$} & & \\

\midrule

\multirow{2}{*}{A (growth)} & \hybrid{} & 69 & 0.275 & 0.401 & 0.343 & 0.569 & 0.646 \\

 & OSNet+ctx & 69 & \textbf{0.826} & \textbf{0.647} & \textbf{0.739} & 0.601 & 0.631 \\

\midrule

\multirow{2}{*}{B (photometric)} & \hybrid{} & 82 & 0.274 & 0.403 & 0.342 & 0.574 & 0.640 \\

 & OSNet+ctx & 82 & \textbf{0.671} & \textbf{0.570} & \textbf{0.663} & 0.594 & 0.609 \\

\bottomrule

\end{tabular}

\end{table}